% ============================================================
\chapter{Syst\`emes Dynamiques --- D\'efinitions et Exemples}
\label{ch:definitions}
% ============================================================

En 1889, Henri Poincar\'e participe \`a un concours organis\'e par le roi Oscar
II de Su\`ede~: d\'emontrer la stabilit\'e du syst\`eme solaire. Non seulement
Poincar\'e ne parvient pas \`a prouver la stabilit\'e, mais il d\'ecouvre
quelque chose de bien plus profond~: l'existence de comportements si
compliqu\'es que toute pr\'ediction \`a long terme devient impossible. C'est
l\`a que na\^it la th\'eorie des syst\`emes dynamiques~: l'\'etude de la
fa\c{c}on dont les \'etats d'un syst\`eme \'evoluent au fil du temps, avec
toute la richesse --- ordre, chaos, bifurcations --- que cela implique.

% ------------------------------------------------------------
\section{Introduction}
% ------------------------------------------------------------

Un syst\`eme dynamique est un mod\`ele math\'ematique d\'ecrivant l'\'evolution
temporelle d'un \'etat dans un espace donn\'e. Cette \'evolution peut \^etre
r\'egie par une \'equation diff\'erentielle (temps continu) ou par une
application it\'er\'ee (temps discret). Dans ce chapitre, nous posons les
d\'efinitions fondamentales et illustrons ces concepts par des exemples
classiques issus de la physique, de la biologie et de l'ing\'enierie.

\section{Syst\`emes dynamiques continus}

\begin{definition}[Syst\`eme dynamique autonome continu]
Un \textbf{syst\`eme dynamique autonome continu} sur un ouvert
$U \subseteq \R^n$ est d\'efini par l'\'equation diff\'erentielle ordinaire
\[
    \dot{x} = f(x), \qquad x \in U,
\]
o\`u $f : U \to \R^n$ est un champ de vecteurs au moins de classe $C^1$.
\end{definition}

\begin{definition}[Syst\`eme non autonome]
Un \textbf{syst\`eme dynamique non autonome} est de la forme
\[
    \dot{x} = f(t, x), \qquad (t, x) \in \R \times U,
\]
o\`u $f$ d\'epend explicitement du temps $t$.
\end{definition}

\begin{remarque}
Tout syst\`eme non autonome $\dot{x} = f(t,x)$ sur $\R^n$ peut \^etre
transform\'e en un syst\`eme autonome sur $\R^{n+1}$ en posant
$\dot{t} = 1$, ce qui donne le syst\`eme augment\'e :
\[
    \begin{pmatrix} \dot{x} \\ \dot{t} \end{pmatrix} =
    \begin{pmatrix} f(t,x) \\ 1 \end{pmatrix}.
\]
\end{remarque}

\subsection{Existence et unicit\'e des solutions}

\begin{theoreme}[Cauchy-Lipschitz / Picard-Lindel\"of]
\label{thm:cauchy-lipschitz}
Soit $f : U \to \R^n$ localement lipschitzienne sur l'ouvert $U \subseteq \R^n$.
Pour tout $x_0 \in U$, il existe un intervalle maximal $I(x_0) = (\alpha, \omega)$
avec $-\infty \leq \alpha < 0 < \omega \leq +\infty$ et une unique solution
maximale $\varphi(\cdot, x_0) : I(x_0) \to U$ du probl\`eme de Cauchy
\[
    \dot{x} = f(x), \qquad x(0) = x_0.
\]
\end{theoreme}

\begin{proof}
La d\'emonstration classique utilise le th\'eor\`eme du point fixe de Banach
appliqu\'e \`a l'op\'erateur de Picard
\[
    (T\varphi)(t) = x_0 + \int_0^t f(\varphi(s))\, ds
\]
dans l'espace $C([0,h], \overline{B}(x_0, r))$ muni de la norme uniforme,
pour $h > 0$ suffisamment petit. La condition de Lipschitz
$\norm{f(x) - f(y)} \leq L\norm{x - y}$ assure que $T$ est une contraction
pour $h < 1/L$. L'unicit\'e sur l'intervalle maximal d\'ecoule du th\'eor\`eme
de prolongement et du lemme de Gr\"onwall.
\end{proof}

\begin{theoreme}[Comportement aux bornes de l'intervalle maximal]
\label{thm:explosion}
Soit $\varphi(\cdot, x_0) : (\alpha, \omega) \to U$ la solution maximale.
Si $\omega < +\infty$, alors pour tout compact $K \subseteq U$, il existe
$t_K < \omega$ tel que $\varphi(t, x_0) \notin K$ pour tout $t > t_K$.
Autrement dit, la trajectoire \og explose \fg ou quitte tout compact en
temps fini.
\end{theoreme}

\section{Flots et semi-flots}

\begin{definition}[Flot]
Le \textbf{flot} associ\'e au syst\`eme $\dot{x} = f(x)$ est l'application
\[
    \varphi : \mathcal{D} \subseteq \R \times U \to U, \qquad
    (t, x_0) \mapsto \varphi(t, x_0) = \varphi_t(x_0),
\]
o\`u $\mathcal{D}$ est l'ouvert maximal de d\'efinition. Le flot satisfait :
\begin{enumerate}
    \item $\varphi_0 = \mathrm{Id}_U$ (condition initiale) ;
    \item $\varphi_{t+s} = \varphi_t \circ \varphi_s$ pour tous $s, t$
    admissibles (propri\'et\'e de groupe).
\end{enumerate}
\end{definition}

\begin{proposition}[R\'egularit\'e du flot]
Si $f \in C^k(U, \R^n)$ avec $k \geq 1$, alors le flot
$\varphi : \mathcal{D} \to U$ est de classe $C^k$ en $(t, x_0)$.
\end{proposition}

\begin{definition}[Semi-flot]
Un \textbf{semi-flot} est d\'efini pour $t \geq 0$ uniquement. Il satisfait
les m\^emes propri\'et\'es que le flot, mais seulement pour $t, s \geq 0$.
Les \'equations aux d\'eriv\'ees partielles de type parabolique engendrent
naturellement des semi-flots.
\end{definition}

\section{Orbites, trajectoires et ensembles limites}

\begin{definition}[Orbite]
L'\textbf{orbite} (ou \textbf{trajectoire}) passant par $x_0$ est l'ensemble
\[
    \gamma(x_0) = \{\varphi(t, x_0) : t \in I(x_0)\} \subseteq U.
\]
L'\textbf{orbite positive} est $\gamma^+(x_0) = \{\varphi(t, x_0) : t \geq 0\}$
et l'\textbf{orbite n\'egative} est $\gamma^-(x_0) = \{\varphi(t, x_0) : t \leq 0\}$.
\end{definition}

\begin{definition}[Points d'\'equilibre]
Un point $x^* \in U$ est un \textbf{point d'\'equilibre} (ou \textbf{point fixe},
\textbf{point stationnaire}) du syst\`eme $\dot{x} = f(x)$ si $f(x^*) = 0$.
De mani\`ere \'equivalente, $\varphi(t, x^*) = x^*$ pour tout $t$.
\end{definition}

\begin{definition}[Orbite p\'eriodique]
Une orbite $\gamma(x_0)$ est \textbf{p\'eriodique} s'il existe $T > 0$ tel que
$\varphi(T, x_0) = x_0$ et $\varphi(t, x_0) \neq x_0$ pour $0 < t < T$.
Le plus petit tel $T$ est appel\'e la \textbf{p\'eriode minimale}.
\end{definition}

\begin{definition}[Ensemble $\omega$-limite]
L'\textbf{ensemble $\omega$-limite} de $x_0$ est
\[
    \omega(x_0) = \bigcap_{T \geq 0} \overline{\{\varphi(t, x_0) : t \geq T\}}.
\]
Un point $y$ appartient \`a $\omega(x_0)$ si et seulement s'il existe une suite
$t_n \to +\infty$ telle que $\varphi(t_n, x_0) \to y$.
\end{definition}

\begin{definition}[Ensemble $\alpha$-limite]
De mani\`ere analogue, l'\textbf{ensemble $\alpha$-limite} est
\[
    \alpha(x_0) = \bigcap_{T \leq 0} \overline{\{\varphi(t, x_0) : t \leq T\}}.
\]
\end{definition}

\begin{proposition}[Propri\'et\'es des ensembles limites]
\label{prop:omega-props}
Soit $\gamma^+(x_0)$ born\'ee. Alors :
\begin{enumerate}
    \item $\omega(x_0)$ est non vide, compact et connexe ;
    \item $\omega(x_0)$ est invariant par le flot : $\varphi_t(\omega(x_0)) = \omega(x_0)$ pour tout $t$ ;
    \item $\mathrm{dist}(\varphi(t, x_0), \omega(x_0)) \to 0$ quand $t \to +\infty$.
\end{enumerate}
\end{proposition}

\section{Espace des phases et portrait de phase}

\begin{definition}[Portrait de phase]
Le \textbf{portrait de phase} d'un syst\`eme dynamique est la partition de
l'espace des phases $U$ en orbites. Il fournit une description qualitative
globale de la dynamique.
\end{definition}

\begin{remarque}
Deux orbites distinctes ne peuvent pas se croiser (cons\'equence de l'unicit\'e
des solutions). L'espace des phases est donc partitionn\'e en orbites
disjointes.
\end{remarque}

\begin{exemple}[Oscillateur harmonique]
Consid\'erons l'oscillateur harmonique simple
\[
    \ddot{q} + \omega^2 q = 0,
\]
que l'on \'ecrit comme syst\`eme du premier ordre avec $x_1 = q$, $x_2 = \dot{q}$ :
\[
    \begin{cases}
        \dot{x}_1 = x_2, \\
        \dot{x}_2 = -\omega^2 x_1.
    \end{cases}
\]
La matrice du syst\`eme est $A = \begin{pmatrix} 0 & 1 \\ -\omega^2 & 0 \end{pmatrix}$,
de valeurs propres $\lambda = \pm i\omega$. Les orbites sont des ellipses
centr\'ees \`a l'origine. Le portrait de phase est un \textbf{centre}.
\end{exemple}

\begin{exemple}[Pendule simple]
\label{ex:pendule}
Le pendule simple non amorti satisfait
\[
    \ddot{\theta} + \frac{g}{l}\sin\theta = 0.
\]
En posant $x_1 = \theta$, $x_2 = \dot{\theta}$ :
\[
    \begin{cases}
        \dot{x}_1 = x_2, \\
        \dot{x}_2 = -\frac{g}{l}\sin x_1.
    \end{cases}
\]
Les points d'\'equilibre sont $(\theta^*, 0)$ avec $\theta^* = k\pi$, $k \in \Z$.
Les points $\theta^* = 2k\pi$ sont des centres (positions basses) et les points
$\theta^* = (2k+1)\pi$ sont des cols (positions hautes).
\end{exemple}

\begin{intuition}
Le portrait de phase du pendule r\'ev\`ele deux types de mouvement :
\begin{itemize}
    \item Les trajectoires ferm\'ees pr\`es des centres correspondent aux
    oscillations (lib\'erations) ;
    \item Les trajectoires ouvertes loin des centres correspondent aux
    rotations compl\`etes du pendule.
\end{itemize}
Les trajectoires s\'eparant ces deux r\'egions sont les \textbf{s\'eparatrices},
qui relient les points cols.
\end{intuition}

% --- Portrait de phase du pendule ---
\begin{center}
\begin{tikzpicture}
    \begin{axis}[
        width=12cm, height=7cm,
        xlabel={$\theta$}, ylabel={$\dot{\theta}$},
        xmin=-7, xmax=7, ymin=-4, ymax=4,
        axis lines=middle,
        xtick={-6.283,-3.14159,0,3.14159,6.283},
        xticklabels={$-2\pi$,$-\pi$,$0$,$\pi$,$2\pi$},
        title={Portrait de phase du pendule simple}
    ]
    % Separatrices
    \addplot[thick, red, domain=-3.14:3.14, samples=100]
        {sqrt(2*(1+cos(deg(x))))};
    \addplot[thick, red, domain=-3.14:3.14, samples=100]
        {-sqrt(2*(1+cos(deg(x))))};
    % Small oscillations
    \addplot[blue, domain=-2:2, samples=80]
        {sqrt(0.5*(1-cos(deg(2*x)))/2 + 0.3)};
    \addplot[blue, domain=-2:2, samples=80]
        {-sqrt(0.5*(1-cos(deg(2*x)))/2 + 0.3)};
    \end{axis}
\end{tikzpicture}
\end{center}

\section{Syst\`emes dynamiques discrets}

\begin{definition}[Syst\`eme dynamique discret]
Un \textbf{syst\`eme dynamique discret} est d\'efini par l'it\'eration d'une
application $f : U \to U$ :
\[
    x_{n+1} = f(x_n), \qquad n \in \N.
\]
L'orbite de $x_0$ est la suite $(x_0, f(x_0), f^2(x_0), \ldots)$
o\`u $f^n = \underbrace{f \circ f \circ \cdots \circ f}_{n \text{ fois}}$.
\end{definition}

\begin{definition}[Points fixes et p\'eriodiques]
Un point $x^*$ est un \textbf{point fixe} de $f$ si $f(x^*) = x^*$.
Un point $x^*$ est \textbf{p\'eriodique de p\'eriode} $p$ si $f^p(x^*) = x^*$
et $f^k(x^*) \neq x^*$ pour $1 \leq k < p$.
\end{definition}

\begin{exemple}[Application logistique]
L'\textbf{application logistique} est d\'efinie par
\[
    f_r(x) = rx(1-x), \qquad x \in [0,1], \quad r \in [0,4].
\]
Les points fixes satisfont $x = rx(1-x)$, soit $x^* = 0$ et
$x^* = 1 - 1/r$ (pour $r > 1$).
\end{exemple}

\section{Premi\`eres int\'egrales et syst\`emes conservatifs}

\begin{definition}[Premi\`ere int\'egrale]
Une fonction $H : U \to \R$ de classe $C^1$ est une \textbf{premi\`ere
int\'egrale} (ou \textbf{constante du mouvement}) du syst\`eme $\dot{x} = f(x)$
si $H$ est constante le long de toute orbite :
\[
    \frac{d}{dt} H(\varphi(t, x_0)) = \nabla H(\varphi(t, x_0)) \cdot
    f(\varphi(t, x_0)) = 0 \quad \forall t.
\]
\end{definition}

\begin{theoreme}[R\'eduction dimensionnelle]
Si $H$ est une premi\`ere int\'egrale telle que $\nabla H(x) \neq 0$ sur un
ensemble de niveau $\Sigma_c = \{x \in U : H(x) = c\}$, alors $\Sigma_c$
est une sous-vari\'et\'e de dimension $n-1$ invariante par le flot.
La dynamique se r\'eduit donc \`a l'\'etude d'un syst\`eme sur $\Sigma_c$.
\end{theoreme}

\begin{definition}[Syst\`eme hamiltonien]
Un \textbf{syst\`eme hamiltonien} sur $\R^{2n}$ est d\'efini par
\[
    \dot{q}_i = \frac{\partial H}{\partial p_i}, \qquad
    \dot{p}_i = -\frac{\partial H}{\partial q_i}, \qquad i = 1, \ldots, n,
\]
o\`u $H : \R^{2n} \to \R$ est la \textbf{fonction de Hamilton} (\'energie).
Le hamiltonien $H$ est automatiquement une premi\`ere int\'egrale.
\end{definition}

\begin{exemple}[Pendule comme syst\`eme hamiltonien]
Le pendule simple (Exemple~\ref{ex:pendule}) est hamiltonien avec
\[
    H(\theta, p) = \frac{p^2}{2} - \frac{g}{l}\cos\theta,
\]
o\`u $p = \dot{\theta}$. Les courbes de niveau de $H$ co\"incident avec
les orbites du portrait de phase.
\end{exemple}

\section{Syst\`emes dissipatifs et volume de Liouville}

\begin{theoreme}[Formule de Liouville]
\label{thm:liouville}
Soit $\Omega(t) = \varphi_t(\Omega_0)$ l'image d'un domaine mesurable
$\Omega_0 \subseteq U$ par le flot. Alors
\[
    \frac{d}{dt}\mathrm{Vol}(\Omega(t)) = \int_{\Omega(t)} \mathrm{div}\, f(x)\, dx.
\]
\end{theoreme}

\begin{corollaire}
Si $\mathrm{div}\, f(x) = 0$ pour tout $x \in U$, le flot pr\'eserve le volume
(syst\`eme \textbf{incompressible}). C'est le cas de tout syst\`eme hamiltonien.
\end{corollaire}

\begin{definition}[Syst\`eme dissipatif]
Un syst\`eme est \textbf{dissipatif} si $\mathrm{div}\, f(x) < 0$ sur $U$.
Dans ce cas, les volumes dans l'espace des phases contractent au cours du temps.
\end{definition}

\begin{exemple}[Oscillateur amorti]
L'oscillateur harmonique amorti
\[
    \dot{x}_1 = x_2, \qquad \dot{x}_2 = -\omega^2 x_1 - 2\gamma x_2, \quad \gamma > 0,
\]
a pour divergence $\mathrm{div}\, f = -2\gamma < 0$. C'est un syst\`eme dissipatif.
Les volumes contractent au taux $e^{-2\gamma t}$.
\end{exemple}

\section{Exemples de syst\`emes dynamiques c\'el\`ebres}

\begin{exemple}[Syst\`eme de Lotka-Volterra]
Le mod\`ele proie-pr\'edateur de Lotka-Volterra est
\[
    \dot{x} = \alpha x - \beta xy, \qquad \dot{y} = \delta xy - \gamma y,
\]
avec $\alpha, \beta, \gamma, \delta > 0$. Le point d'\'equilibre non trivial est
$(\gamma/\delta, \alpha/\beta)$. Ce syst\`eme poss\`ede la premi\`ere int\'egrale
$H(x,y) = \delta x - \gamma \ln x + \beta y - \alpha \ln y$.
\end{exemple}

\begin{exemple}[Syst\`eme de Lorenz]
Le syst\`eme de Lorenz, introduit en 1963, est
\[
    \dot{x} = \sigma(y - x), \qquad
    \dot{y} = \rho x - y - xz, \qquad
    \dot{z} = xy - \beta z,
\]
avec les param\`etres classiques $\sigma = 10$, $\rho = 28$, $\beta = 8/3$.
La divergence est $\mathrm{div}\, f = -\sigma - 1 - \beta < 0$ :
le syst\`eme est dissipatif.
\end{exemple}

\begin{codeblock}[title=\textbf{Simulation du syst\`eme de Lorenz}]
\begin{minted}{python}
import numpy as np
from scipy.integrate import solve_ivp
import matplotlib.pyplot as plt

def lorenz(t, state, sigma=10, rho=28, beta=8/3):
    x, y, z = state
    return [sigma*(y - x), rho*x - y - x*z, x*y - beta*z]

sol = solve_ivp(lorenz, [0, 50], [1, 1, 1], max_step=0.01,
                dense_output=True)
t = np.linspace(0, 50, 10000)
xyz = sol.sol(t)

fig = plt.figure(figsize=(10, 7))
ax = fig.add_subplot(111, projection='3d')
ax.plot(xyz[0], xyz[1], xyz[2], lw=0.5, color='steelblue')
ax.set_xlabel('x'); ax.set_ylabel('y'); ax.set_zlabel('z')
ax.set_title("Attracteur de Lorenz")
plt.tight_layout()
plt.savefig("lorenz_attractor.pdf")
plt.show()
\end{minted}
\end{codeblock}

\section{\'Equivalence topologique et conjugaison}

\begin{definition}[\'Equivalence topologique]
Deux syst\`emes dynamiques $\dot{x} = f(x)$ et $\dot{y} = g(y)$ sont
\textbf{topologiquement \'equivalents} s'il existe un hom\'eomorphisme
$h : U \to V$ qui envoie les orbites de $f$ sur les orbites de $g$
en pr\'eservant l'orientation (le sens de parcours).
\end{definition}

\begin{definition}[Conjugaison topologique]
Deux syst\`emes sont \textbf{topologiquement conjugu\'es} si le
changement de variable $h$ respecte aussi la param\'etrisation temporelle :
$h \circ \varphi_t^f = \varphi_t^g \circ h$ pour tout $t$.
\end{definition}

\begin{remarque}
La conjugaison topologique est plus forte que l'\'equivalence topologique.
La conjugaison pr\'eserve les p\'eriodes des orbites p\'eriodiques,
contrairement \`a la simple \'equivalence.
\end{remarque}

\section{Exercices}

\begin{exercice}[Points d'\'equilibre]
Trouver tous les points d'\'equilibre du syst\`eme
\[
    \dot{x} = y - x^2, \qquad \dot{y} = x - 2.
\]
D\'eterminer la nature de chaque point d'\'equilibre en calculant la matrice
jacobienne et ses valeurs propres.
\end{exercice}

\begin{exercice}[Premi\`ere int\'egrale]
Montrer que $H(x, y) = x^2 + y^2 - 2\ln(1 + x^2 + y^2)$ est une premi\`ere
int\'egrale du syst\`eme
\[
    \dot{x} = -y\,\frac{x^2 + y^2}{1 + x^2 + y^2}, \qquad
    \dot{y} = x\,\frac{x^2 + y^2}{1 + x^2 + y^2}.
\]
\end{exercice}

\begin{exercice}[Flot explicite]
Calculer explicitement le flot du syst\`eme lin\'eaire
$\dot{x} = Ax$ avec $A = \begin{pmatrix} 0 & 1 \\ -1 & 0 \end{pmatrix}$.
V\'erifier la propri\'et\'e de groupe $\varphi_{t+s} = \varphi_t \circ \varphi_s$.
\end{exercice}

\begin{exercice}[Syst\`eme dissipatif]
Calculer la divergence du syst\`eme de R\"ossler
\[
    \dot{x} = -y - z, \qquad \dot{y} = x + ay, \qquad \dot{z} = b + z(x - c),
\]
et d\'eterminer pour quelles valeurs des param\`etres $a, b, c$ le syst\`eme
est dissipatif.
\end{exercice}

\begin{exercice}[Application logistique]
Pour l'application logistique $f_r(x) = rx(1-x)$ :
\begin{enumerate}
    \item Trouver les points fixes et \'etudier leur stabilit\'e en fonction de $r$.
    \item Trouver les orbites de p\'eriode~2 et \'etudier leur stabilit\'e.
    \item Montrer que pour $r > 4$, l'orbite de presque tout point initial
    dans $[0,1]$ s'\'echappe vers $-\infty$.
\end{enumerate}
\end{exercice}

\begin{exercice}[Volume et Liouville]
Soit le syst\`eme $\dot{x} = f(x)$ sur $\R^3$ avec
$f(x,y,z) = (y, -x - \epsilon y, z - \epsilon z^3)$.
\begin{enumerate}
    \item Calculer $\mathrm{div}\, f$.
    \item Pour quelles valeurs de $\epsilon$ le syst\`eme est-il dissipatif ?
    conservatif ?
    \item D\'ecrire qualitativement le comportement des volumes au cours du temps.
\end{enumerate}
\end{exercice}
